%==============================================================================
% CHAPTER 4
%==============================================================================
\chapter{Results and stress-testing}
\markboth{Chapter 4 : Results and stress-testing}{}

This chapter gathers the results obtained : first the physical behaviour that a
reference simulation reproduces, then the quantitative validation of the
surrogate, and finally what happened when we set out to break the agent on
purpose. The last part is the longest, and it is where we learned the most.

\section{Verifying the solver against analytical solutions}
\label{sec:verification}

Before exploiting anything the chain produces, we had to satisfy ourselves that the
solver underneath it is correct. All the results presented later
in this chapter are matters of \emph{internal consistency} : the solver is
consistent with itself, the surrogate is consistent with the solver. None
confronts the chain with an external truth.

For this purpose OFFBEAT distributes a suite of verification cases, some of
which have a \textbf{closed-form analytical solution}. We ran three of them, covering three
distinct areas of physics, on our own installation and compared the output with
the reference solution (table~\ref{tab:verif}).

\begin{table}[htbp]
\centering
\caption[Verification against analytical solutions]{Verification of the solver
against analytical solutions, run in the internship environment (OpenFOAM
v2506, OFFBEAT \textit{master} branch). The indicators are those defined by the
verification scripts shipped with the code.}
\label{tab:verif}
\small
\begin{tabularx}{\textwidth}{@{}>{\raggedright\arraybackslash}p{3.4cm}Xl@{}}
\toprule
\textbf{Verification case} & \textbf{Physics exercised} &
\textbf{Deviation from analytical}\\
\midrule
Thick cylinder expansion &
  Elastoplastic mechanics; radial stress at the inner wall,
  \num{150} points compared &
  \SI{1.02}{\percent} (max)\newline \SI{0.82}{\percent} (mean)\\
\addlinespace
Porosity redistribution &
  Porosity migration under a thermal gradient &
  \SI{0.064}{\percent} ($t=\SI{180}{\second}$)\newline
  \SI{0.16}{\percent} ($t=\SI{360}{\second}$)\\
\addlinespace
Actinide redistribution &
  Thermal diffusion of actinides, Clément's analytical solution &
  \num{5.4e-3}~\% (Pu)\newline \num{3.2e-3}~\% (Am)\\
\bottomrule
\end{tabularx}
\end{table}

Agreement is good on all three cases : at worst one per cent on the mechanical
case, which is also the most demanding since it involves a plastification front,
and deviations of the order of a thousandth of a per cent on the two fuel
physics cases.

This result does not establish the validity of the agent; it establishes that
of the solver the agent drives, and therefore that any errors observed
downstream are attributable to the tooled chain and not to the computational
core. This is precisely the separation needed to interpret the correctness
defects of section~\ref{sec:bugs} : none of them came from OFFBEAT.

\vspace{0.2cm}
\noindent
\textit{Remark.} A fourth case in the suite (laser heating) could not be run : it
fails while reading its own dictionaries under OpenFOAM v2506. This is an
incompatibility between that case and that version of OpenFOAM, independent of
the work presented here, but it illustrates how fragile verification suites are
across version upgrades.

\section{Reference simulation over an irradiation cycle}

\subsection{Configuration}

The reference case is a PWR rod subjected to a linear heat rate of
\SI{25}{\kilo\watt\per\metre} over a simulated duration of \SI{6.3e7}{\second},
about two years. The geometry is that of the reference template : pellet radius
\SI{4.5}{\milli\metre}, cladding inner radius \SI{4.565}{\milli\metre}, hence an
initial gap of \SI{65}{\micro\metre}. Seven states were written during the run.

\subsection{Thermal evolution}

Figure~\ref{fig:temp} shows the evolution of the maximum temperatures of the
pellet and of the cladding.

\begin{figure}[htbp]
\centering
\begin{tikzpicture}
\begin{axis}[
  width=0.86\textwidth, height=6.2cm,
  xlabel={Irradiation time [days]},
  ylabel={Maximum temperature [K]},
  xmin=0, xmax=760, ymin=250, ymax=1600,
  grid=major, grid style={cGrid!60,very thin},
  legend pos=north east, legend style={font=\scriptsize,draw=black!30},
  label style={font=\small}, tick label style={font=\scriptsize},
  every axis plot/.append style={thick,mark size=1.6pt},
]
\addplot[color=cFuel,mark=*] coordinates {
  (170.0,1361.75) (354.1,1385.75) (456.9,1488.66)
  (590.2,1444.20) (710.9,1461.03) (729.2,1467.29)};
\addlegendentry{Pellet (centreline)}
\addplot[color=cClad,mark=square*] coordinates {
  (0,300.0) (170.0,639.09) (354.1,637.71) (456.9,641.51)
  (590.2,637.90) (710.9,637.39) (729.2,637.39)};
\addlegendentry{Cladding (PCT)}
\end{axis}
\end{tikzpicture}
\caption[Pellet and cladding temperatures over the cycle]{Evolution of the
maximum pellet and cladding temperatures over the cycle. Both are maxima taken
over the whole rod, which is why the pellet curve is not monotonic : the
hottest point moves as the axial power profile shifts.}
\label{fig:temp}
\end{figure}

The cladding temperature remains very stable, around \SI{637}{\kelvin},
which is consistent because it is essentially imposed by the coolant.

The pellet curve on the contrary does not rise monotonically, it peaks at \SI{1489}{\kelvin} on day
\num{457} and falls back to \SI{1444}{\kelvin}, then rises again. Degradation of
the fuel conductivity cannot explain that, because it only goes one way.

Our first explanation was that the curve is a maximum taken over the whole rod,
so it merely follows the hottest point as that point moves. Checking cell by
cell showed this to be incomplete. The temperature is non-monotonic at
\emph{every} axial position, not only where the maximum happens to sit : at
\SI{0.15}{\metre} it goes \SI{1150}{\kelvin}, \SI{1242}{\kelvin},
\SI{1117}{\kelvin}, then back up. Something makes individual points cool down.

That something is the axial power profile, which the template tabulates at
twelve instants and which cycles between three shapes. It does not simply move
a peak : it \textbf{redistributes the power along the rod}, so the local power
density at a fixed height is not constant either. Over the cycle it varies by
\SI{14}{\percent} at mid-height and by up to \SI{49}{\percent} near the ends.
A point that receives less power cools down, however degraded the fuel around it
has become.

Temperature is therefore the wrong quantity in which to look for the
conductivity. The right one is the \textbf{thermal resistance} of the pellet,
the centre-to-edge drop divided by the power density that produces it. That
ratio rises by \SIrange{24}{27}{\percent} at every height, and monotonically at
seven of the ten positions sampled ; the three exceptions are the rod ends,
where the local power swings so widely that the extraction becomes noisy. Since
this ratio is proportional to $1/k$, it puts the loss of conductivity over the
cycle at about \SI{20}{\percent}.

The picture is then consistent : writing
$T_{\text{centre}} \simeq T_{\text{edge}} + R_{\text{th}}\,q$, the resistance
$R_{\text{th}}$ rises monotonically as the fuel degrades, while the local power
$q$ oscillates as the profile reshapes. The temperature inherits the
oscillation, which is why no curve of $T$ can be monotonic.

One detail confirms the picture from the other side : over the same period the
pellet \emph{outer} temperature falls, from \SI{712}{\kelvin} to
\SI{673}{\kelvin}. The pellet is hotter at heart and cooler at its edge because
the gap is closing and evacuating heat better, which is the same closure that
figure~\ref{fig:gap-stress} dates independently.

\subsection{Gap closure and stress reversal}

The most interesting result appears when the evolution of the pellet-cladding
gap is set against that of the hoop stress (figure~\ref{fig:gap-stress}).

\begin{figure}[htbp]
\centering
\begin{minipage}{0.48\textwidth}
\centering
\begin{tikzpicture}
\begin{axis}[
  width=\textwidth, height=5.4cm,
  xlabel={Time [days]}, ylabel={Gap width [\si{\micro\metre}]},
  xmin=0, xmax=760, ymin=-11, ymax=13,
  grid=major, grid style={cGrid!60,very thin},
  label style={font=\scriptsize}, tick label style={font=\scriptsize},
  title style={font=\scriptsize\bfseries}, title={(a) Pellet-cladding gap},
]
\addplot[color=cAlert,thick,mark=*,mark size=1.6pt] coordinates {
  (170.0,10.097) (354.1,0.370) (456.9,-1.959)
  (590.2,-5.251) (710.9,-7.703) (729.2,-7.997)};
\addplot[black!55,dashed,thick,domain=0:760,forget plot] {0};
\node[font=\tiny,anchor=west,text=black!65] at (axis cs:250,1.8) {contact};
\end{axis}
\end{tikzpicture}
\end{minipage}\hfill
\begin{minipage}{0.48\textwidth}
\centering
\begin{tikzpicture}
\begin{axis}[
  width=\textwidth, height=5.4cm,
  xlabel={Time [days]},
  ylabel={Hoop stress [MPa]},
  xmin=0, xmax=760, ymin=-60, ymax=100,
  grid=major, grid style={cGrid!60,very thin},
  label style={font=\scriptsize}, tick label style={font=\scriptsize},
  title style={font=\scriptsize\bfseries},
  title={(b) Stress $\sigma_{\theta\theta}$ (cladding)},
]
\addplot[color=cClad,thick,mark=square*,mark size=1.6pt] coordinates {
  (170.0,-43.01) (354.1,-38.23) (456.9,-31.33)
  (590.2,28.79) (710.9,71.59) (729.2,79.11)};
\addplot[black!55,dashed,thick,domain=0:760,forget plot] {0};
\node[font=\tiny,anchor=east,text=black!65] at (axis cs:740,-16) {compression};
\node[font=\tiny,anchor=east,text=black!65] at (axis cs:740,16) {tension};
\end{axis}
\end{tikzpicture}
\end{minipage}
\caption[Gap closure and hoop stress]{Pellet-cladding gap closure (a) and hoop
stress in the cladding (b). The change of sign of the stress coincides with the
establishment of contact : the cladding, first compressed by the coolant, is then
put into tension by the pellet.}
\label{fig:gap-stress}
\end{figure}

These two curves are the main argument in favor of the model. As long as the gap remains open, the cladding is
\textbf{compressed} by the coolant pressure : the stress is negative (about\SI{-43}{\mega\pascal}). The gap closes between days 
354 and 457. The stress then becomes \textbf{positive} and grows to
\SI{79}{\mega\pascal}. This is the signature of PCMI : the pellet, as it
swells, puts the cladding into tension. The automatic prognosis (block~D2)
places the crossing of the closure threshold at about \SI{2.33e7}{\second}, that
is at one third of the cycle, consistent with the curve.

\subsection{Spatial distribution of the temperature}

Figure~\ref{fig:carte} maps the temperature over the whole rod at end of cycle.
The mesh is a \num{40}$\times$\num{11} grid in $(r,z)$. The radial scale is exaggerated
by a factor of \num{600}, the rod is \SI{5.3}{\milli\metre} in radius for
\SI{3.2}{\metre} in height. If it was to scale, it would be a impossible to see.
